\documentclass[11pt,a4paper,openany]{article}
% ==============================================================================
% ESTILO_FISCAL.TEX — CONFIGURACIÓN VISUAL PARA DERECHO FISCAL (UGR / repasaYA)
% ==============================================================================

\usepackage[utf8]{inputenc}
\usepackage[T1]{fontenc}
\usepackage[spanish,es-tabla]{babel}
\usepackage{lmodern}
\usepackage{microtype}
\usepackage[dvipsnames,svgnames,table,xcdraw]{xcolor}
\usepackage{amsmath,amssymb,amsfonts}
\usepackage{graphicx}
\usepackage{booktabs}
\usepackage{tabularx}
\usepackage{longtable}
\usepackage{array}
\usepackage{eurosym}
\usepackage{enumitem}
\usepackage{fancyhdr}
\usepackage{titlesec}
\usepackage[breakall]{truncate}
\usepackage{tcolorbox}
\tcbuselibrary{skins,breakable}

% Configuración de Márgenes
\usepackage[a4paper,top=2.5cm,bottom=2.5cm,left=2.5cm,right=2.5cm,headheight=24pt]{geometry}

% Paleta de colores sobria, elegante y jurídica
\definecolor{fiscalDark}{RGB}{15, 30, 54}
\definecolor{fiscalPrimary}{RGB}{24, 60, 108}
\definecolor{fiscalSecondary}{RGB}{41, 98, 153}
\definecolor{fiscalAccent}{RGB}{180, 83, 9}
\definecolor{fiscalLight}{RGB}{248, 250, 252}
\definecolor{fiscalBorder}{RGB}{203, 213, 225}
\definecolor{fiscalGreen}{RGB}{22, 101, 52}
\definecolor{fiscalGreenLight}{RGB}{240, 253, 244}
\definecolor{fiscalRed}{RGB}{153, 27, 27}
\definecolor{fiscalRedLight}{RGB}{254, 242, 242}
\definecolor{fiscalYellowLight}{RGB}{254, 252, 232}

% Hipervínculos
\usepackage{hyperref}
\hypersetup{
    colorlinks=true,
    linkcolor=fiscalPrimary,
    citecolor=fiscalPrimary,
    urlcolor=fiscalSecondary,
    pdfauthor={repasaYA - Facultad de CC.EE. y Empresariales UGR},
    pdftitle={Derecho Fiscal - Material de Examen},
    pdfsubject={IRPF, LGT, Casos Prácticos Oficiales y Exámenes Resueltos UGR}
}

% Encabezados y Pies de página seguros (sin colisiones)
\pagestyle{fancy}
\fancyhf{}
\renewcommand{\headrulewidth}{0.5pt}
\renewcommand{\footrulewidth}{0.4pt}
\fancyhead[LE]{\footnotesize\bfseries\textsf{\thepage}\hspace{1.2em}\textsf{\color{fiscalPrimary}\truncate{0.75\textwidth}{\leftmark}}}
\fancyhead[RO]{\footnotesize\textsf{\color{fiscalPrimary}\truncate{0.75\textwidth}{\rightmark}}\hspace{1.2em}\bfseries\textsf{\thepage}}
\fancyfoot[C]{\footnotesize\textsf{\color{gray}repasaYA $\cdot$ Derecho Fiscal (5º ADE / DADE) $\cdot$ Facultad de CC.EE. y Empresariales (UGR)}}

% Estilos de Secciones
\titleformat{\section}
  {\Large\bfseries\color{fiscalDark}\sffamily}
  {\thesection}{1em}{}[\titlerule]

\titleformat{\subsection}
  {\large\bfseries\color{fiscalPrimary}\sffamily}
  {\thesubsection}{0.8em}{}

\titleformat{\subsubsection}
  {\normalsize\bfseries\color{fiscalSecondary}\sffamily}
  {\thesubsubsection}{0.6em}{}

% Cajas especializadas tcolorbox
\newtcolorbox{enunciadoBox}[1][]{
    enhanced,
    breakable,
    colback=fiscalLight,
    colframe=fiscalPrimary,
    arc=3pt,
    boxrule=1.2pt,
    title={\textbf{\textsf{#1}}},
    coltitle=white,
    attach boxed title to top left={yshift=-2mm, xshift=4mm},
    boxed title style={colback=fiscalPrimary, arc=2pt, boxrule=0pt},
    top=4mm, bottom=3mm, left=3mm, right=3mm
}

\newtcolorbox{solucionBox}[1][]{
    enhanced,
    breakable,
    colback=white,
    colframe=fiscalSecondary!70,
    arc=3pt,
    boxrule=0.8pt,
    title={\textbf{\textsf{\color{fiscalDark}#1}}},
    coltitle=fiscalDark,
    attach boxed title to top left={yshift=-2mm, xshift=4mm},
    boxed title style={colback=fiscalLight, arc=2pt, boxrule=0.5pt, colframe=fiscalSecondary!40},
    top=4mm, bottom=3mm, left=3mm, right=3mm
}

\newtcolorbox{resultadoBox}[1][]{
    enhanced,
    colback=fiscalGreenLight,
    colframe=fiscalGreen,
    arc=3pt,
    boxrule=1pt,
    title={\textbf{\textsf{\color{fiscalGreen!80!black}#1}}},
    coltitle=fiscalGreen!80!black,
    attach boxed title to top left={yshift=-2mm, xshift=4mm},
    boxed title style={colback=fiscalGreenLight, arc=2pt, boxrule=0.5pt, colframe=fiscalGreen},
    top=3mm, bottom=3mm, left=3mm, right=3mm
}

\newtcolorbox{trampaBox}[1][]{
    enhanced,
    colback=fiscalYellowLight,
    colframe=fiscalAccent,
    arc=3pt,
    boxrule=1pt,
    title={\textbf{\textsf{\color{fiscalAccent!90!black}⚠️ TRAMPA DE EXAMEN / REGLA DE ORO: #1}}},
    coltitle=fiscalAccent!90!black,
    attach boxed title to top left={yshift=-2mm, xshift=4mm},
    boxed title style={colback=fiscalYellowLight, arc=2pt, boxrule=0.5pt, colframe=fiscalAccent},
    top=3mm, bottom=3mm, left=3mm, right=3mm
}


\begin{document}

% ------------------------------------------------------------------------------
% PORTADA INSTITUCIONAL
% ------------------------------------------------------------------------------
\begin{titlepage}
    \centering
    \vspace*{1.2cm}
    
    {\Huge\textbf{\textsf{\color{fiscalPrimary}DERECHO FISCAL}}\par}
    \vspace{0.35cm}
    {\Large\textbf{\textsf{\color{fiscalAccent}COMPILACIÓN OFICIAL DE EXÁMENES RESUELTOS}}\par}
    \vspace{0.5cm}
    {\large\textsf{Facultad de Ciencias Económicas y Empresariales --- Universidad de Granada}\par}
    \vspace{0.25cm}
    {\normalsize\textsf{Grado en ADE $\cdot$ Doble Grado DADE / FICO}\par}
    
    \vspace{1.2cm}
    
    \begin{tcolorbox}[
        enhanced,
        colback=fiscalLight,
        colframe=fiscalPrimary,
        arc=4pt,
        boxrule=1.5pt,
        width=0.96\textwidth,
        halign=center
    ]
        \vspace{0.2cm}
        {\large\bfseries\color{fiscalDark}EXÁMENES Y PARCIALES OFICIALES ÍNTEGROS VERBATIM}\\[0.3cm]
        \normalsize
        \begin{itemize}[leftmargin=1.5cm,itemsep=3pt]
            \item \textbf{Bloque I: Parcial Oficial Wuolah (Modelos A y B)} --- 23 preguntas de Tema 1 (LGT, tributos, liquidación, deudas y sucesores).
            \item \textbf{Bloque II: Simulacro Oficial Parcial 9 Noviembre} --- 10 preguntas teóricas LGT/IRPF + Caso Práctico Dña. Elena (5 preguntas de liquidación).
            \item \textbf{Bloque III: Examen Ordinario 2019 (Profesora Titular)} --- Enunciado oficial íntegro del matrimonio Dña. María y D. Julio + 18 preguntas resueltas.
            \item \textbf{Bloque IV: Examen Ordinario Enero 2020 (Profesora Titular)} --- Enunciado oficial íntegro del caso D. Jacobo y Dña. Dolores + 20 preguntas resueltas.
        \end{itemize}
        \vspace{0.2cm}
        \textbf{\color{fiscalAccent}76 Preguntas Verbatim $\cdot$ Soluciones de Pizarra $\cdot$ Artículos LGT/LIRPF $\cdot$ Sin Palabrería}
        \vspace{0.2cm}
    \end{tcolorbox}
    
    \vfill
    {\footnotesize\textsf{\color{gray}repasaYA $\cdot$ Material docente de máxima fidelidad para examen $\cdot$ Noviembre 2026}\par}
\end{titlepage}

\newpage
\tableofcontents
\newpage

\section{Bloque I: Parcial Oficial de Nociones Tributarias y LGT (Wuolah)}
\label{sec:wuolah}

\begin{enunciadoBox}[Contexto del Parcial Oficial de Nociones Tributarias]
Examen parcial oficial correspondiente al Bloque Temático I (Concepto de tributo, hechos imponibles, exención/no sujeción, sujetos pasivos, métodos de estimación de bases, extinción de la deuda tributaria y sucesión mortis causa). Incluye los dos modelos oficiales (A y B) con sus plantillas corregidas y la cita exacta de los preceptos de la Ley 58/2003 General Tributaria (LGT).
\end{enunciadoBox}

\vspace{0.4cm}
\subsection{Modelo A: 15 Preguntas Oficiales Verbatim}

\noindent\textbf{\sffamily Pregunta A.1:} En la no sujeción:\\[0.15cm]
\begin{enumerate}[label=\textbf{\alph*)},leftmargin=1.8em,itemsep=1.5pt,topsep=2pt,partopsep=0pt,parsep=0pt]
    \item El sujeto pasivo realiza el hecho imponible de un tributo, pero la ley establece que no quede sujeto al mismo. *(Eso es la exención).*
    \item El sujeto pasivo no realiza el hecho imponible y, sin embargo, queda sujeto a gravamen.
    \item El sujeto pasivo no realiza ningún hecho imponible tipificado y no queda sujeto a gravamen siempre y cuando la ley así lo prevea.
    \item El sujeto pasivo no realiza el hecho imponible y no queda sujeto a gravamen. *(Art. 20.2 LGT: la no sujeción existe porque no se realiza el hecho imponible, aunque la ley no lo mencione expresamente).*
\end{enumerate}
\begin{solucionBox}[Respuesta Oficial: Opción D]
\textbf{Justificación Oficial / Fundamentación:}\\[0.1cm]

\end{solucionBox}
\vspace{0.35cm}

\noindent\textbf{\sffamily Pregunta A.2:} La finalidad del tributo es:\\[0.15cm]
\begin{enumerate}[label=\textbf{\alph*)},leftmargin=1.8em,itemsep=1.5pt,topsep=2pt,partopsep=0pt,parsep=0pt]
    \item Financiar los gastos privados y públicos de las Administraciones Públicas.
    \item Financiar los gastos públicos del Estado.
    \item Financiar los gastos públicos de las Comunidades Autónomas.
    \item Financiar los gastos públicos de las Administraciones Públicas. *(Art. 2.1 LGT).*
\end{enumerate}
\begin{solucionBox}[Respuesta Oficial: Opción D]
\textbf{Justificación Oficial / Fundamentación:}\\[0.1cm]

\end{solucionBox}
\vspace{0.35cm}

\noindent\textbf{\sffamily Pregunta A.3:} Las obligaciones formales que surgen de la relación jurídico-tributaria:\\[0.15cm]
\begin{enumerate}[label=\textbf{\alph*)},leftmargin=1.8em,itemsep=1.5pt,topsep=2pt,partopsep=0pt,parsep=0pt]
    \item Son de carácter pecuniario.
    \item Son pecuniarias y no pecuniarias.
    \item Son no pecuniarias. *(Art. 29.1 LGT).*
    \item Son obligaciones de no hacer o de omisión.
\end{enumerate}
\begin{solucionBox}[Respuesta Oficial: Opción C]
\textbf{Justificación Oficial / Fundamentación:}\\[0.1cm]

\end{solucionBox}
\vspace{0.35cm}

\noindent\textbf{\sffamily Pregunta A.4:} La autoliquidación tributaria:\\[0.15cm]
\begin{enumerate}[label=\textbf{\alph*)},leftmargin=1.8em,itemsep=1.5pt,topsep=2pt,partopsep=0pt,parsep=0pt]
    \item Es un documento en el que la Administración Tributaria procede a determinar el importe de la deuda tributaria.
    \item Es una declaración en la que el sujeto pasivo declara los hechos, califica e interpreta las normas y cuantifica el importe de la deuda tributaria. *(Art. 120.1 LGT).*
    \item Es una declaración solamente de los hechos que tienen incidencia en relación con la deuda tributaria.
    \item Es una declaración de los hechos y una calificación de los mismos, llevada a cabo en todo caso por el contribuyente.
\end{enumerate}
\begin{solucionBox}[Respuesta Oficial: Opción B]
\textbf{Justificación Oficial / Fundamentación:}\\[0.1cm]

\end{solucionBox}
\vspace{0.35cm}

\noindent\textbf{\sffamily Pregunta A.5:} Los impuestos indirectos:\\[0.15cm]
\begin{enumerate}[label=\textbf{\alph*)},leftmargin=1.8em,itemsep=1.5pt,topsep=2pt,partopsep=0pt,parsep=0pt]
    \item Son además impuestos subjetivos.
    \item Además son impuestos periódicos.
    \item Como regla general, con excepciones, permiten la repercusión de la cuota tributaria sobre un tercero. *(Regla general: IVA e II.EE. se repercuten; excepción: ITPAJD).*
    \item Como regla general, con excepciones, no permiten la repercusión de la cuota tributaria sobre un tercero.
\end{enumerate}
\begin{solucionBox}[Respuesta Oficial: Opción C]
\textbf{Justificación Oficial / Fundamentación:}\\[0.1cm]

\end{solucionBox}
\vspace{0.35cm}

\noindent\textbf{\sffamily Pregunta A.6:} El método de estimación objetiva:\\[0.15cm]
\begin{enumerate}[label=\textbf{\alph*)},leftmargin=1.8em,itemsep=1.5pt,topsep=2pt,partopsep=0pt,parsep=0pt]
    \item Es un método obligatorio, pero cabe renuncia.
    \item Es un método obligatorio por imposición de la Ley General Tributaria.
    \item Es un método voluntario. *(Art. 52 LGT: ``tendrá, en todo caso, carácter voluntario para los obligados tributarios'').*
    \item Es un método voluntario u obligatorio, según lo establezca la Ley específica del tributo.
\end{enumerate}
\begin{solucionBox}[Respuesta Oficial: Opción C]
\textbf{Justificación Oficial / Fundamentación:}\\[0.1cm]

\end{solucionBox}
\vspace{0.35cm}

\noindent\textbf{\sffamily Pregunta A.7:} El hecho imponible de una contribución especial puede ser:\\[0.15cm]
\begin{enumerate}[label=\textbf{\alph*)},leftmargin=1.8em,itemsep=1.5pt,topsep=2pt,partopsep=0pt,parsep=0pt]
    \item La realización de negocios, actos o hechos de contenido económico. *(Impuesto).*
    \item La realización de obras públicas y privadas.
    \item La prestación de servicios públicos. *(Tasa).*
    \item El establecimiento de servicios públicos. *(Art. 2.2.b LGT: realización de obras públicas o establecimiento o ampliación de servicios públicos).*
\end{enumerate}
\begin{solucionBox}[Respuesta Oficial: Opción D]
\textbf{Justificación Oficial / Fundamentación:}\\[0.1cm]

\end{solucionBox}
\vspace{0.35cm}

\noindent\textbf{\sffamily Pregunta A.8:} Son impuestos indirectos:\\[0.15cm]
\begin{enumerate}[label=\textbf{\alph*)},leftmargin=1.8em,itemsep=1.5pt,topsep=2pt,partopsep=0pt,parsep=0pt]
    \item El IVA y los Impuestos Especiales.
    \item El Impuesto de Transmisiones Patrimoniales y el Impuesto de Sucesiones y Donaciones. *(El ISD es directo).*
    \item El IVA, los Impuestos Especiales y el Impuesto de Transmisiones Patrimoniales.
    \item Las respuestas a y b son correctas.
\end{enumerate}
\begin{solucionBox}[Respuesta Oficial: Opción C]
\textbf{Justificación Oficial / Fundamentación:}\\[0.1cm]

\end{solucionBox}
\vspace{0.35cm}

\noindent\textbf{\sffamily Pregunta A.9:} La realización del hecho imponible:\\[0.15cm]
\begin{enumerate}[label=\textbf{\alph*)},leftmargin=1.8em,itemsep=1.5pt,topsep=2pt,partopsep=0pt,parsep=0pt]
    \item Origina el nacimiento del pago del tributo.
    \item Origina el nacimiento de la obligación tributaria principal. *(Art. 20.1 LGT).*
    \item Origina el nacimiento de todas las obligaciones tributarias.
    \item a y b son correctas.
\end{enumerate}
\begin{solucionBox}[Respuesta Oficial: Opción B]
\textbf{Justificación Oficial / Fundamentación:}\\[0.1cm]

\end{solucionBox}
\vspace{0.35cm}

\noindent\textbf{\sffamily Pregunta A.10:} Por ocupar parte de la vía pública en un determinado municipio para poner una terraza con mesas durante los meses de verano, el dueño del bar que quiere esa ocupación debe:\\[0.15cm]
\begin{enumerate}[label=\textbf{\alph*)},leftmargin=1.8em,itemsep=1.5pt,topsep=2pt,partopsep=0pt,parsep=0pt]
    \item Pagar un precio.
    \item Pagar un precio público.
    \item Pagar una tasa. *(Art. 2.2.a LGT: utilización privativa o aprovechamiento especial del dominio público).*
    \item Pagar una contribución especial.
\end{enumerate}
\begin{solucionBox}[Respuesta Oficial: Opción C]
\textbf{Justificación Oficial / Fundamentación:}\\[0.1cm]

\end{solucionBox}
\vspace{0.35cm}

\noindent\textbf{\sffamily Pregunta A.11:} El contribuyente:\\[0.15cm]
\begin{enumerate}[label=\textbf{\alph*)},leftmargin=1.8em,itemsep=1.5pt,topsep=2pt,partopsep=0pt,parsep=0pt]
    \item Es el que realiza el hecho imponible. *(Art. 36.2 LGT: ``Es contribuyente el sujeto pasivo que realiza el hecho imponible''. En la opción c se define al sujeto pasivo en general —art. 36.1—, que puede ser contribuyente o sustituto).*
    \item Es el obligado al pago del tributo.
    \item Es el sujeto pasivo obligado a cumplir con las obligaciones materiales y formales derivadas del tributo por imposición de la ley.
    \item Las respuestas a y c son correctas.
\end{enumerate}
\begin{solucionBox}[Respuesta Oficial: Opción A]
\textbf{Justificación Oficial / Fundamentación:}\\[0.1cm]

\end{solucionBox}
\vspace{0.35cm}

\noindent\textbf{\sffamily Pregunta A.12:} El método de estimación directa:\\[0.15cm]
\begin{enumerate}[label=\textbf{\alph*)},leftmargin=1.8em,itemsep=1.5pt,topsep=2pt,partopsep=0pt,parsep=0pt]
    \item Lo aplica la Administración.
    \item Lo aplica el contribuyente.
    \item Constituye junto con el de estimación indirecta, los métodos principales de determinación de bases.
    \item Lo puede aplicar tanto el contribuyente como la Administración. *(Art. 51 LGT).*
\end{enumerate}
\begin{solucionBox}[Respuesta Oficial: Opción D]
\textbf{Justificación Oficial / Fundamentación:}\\[0.1cm]

\end{solucionBox}
\vspace{0.35cm}

\noindent\textbf{\sffamily Pregunta A.13:} Son obligaciones tributarias accesorias:\\[0.15cm]
\begin{enumerate}[label=\textbf{\alph*)},leftmargin=1.8em,itemsep=1.5pt,topsep=2pt,partopsep=0pt,parsep=0pt]
    \item Las sanciones tributarias. *(Nunca).*
    \item La obligación de llevar libros de contabilidad. *(Formal).*
    \item El recargo de apremio. *(Art. 25.1 y 28 LGT).*
    \item La obligación de presentar la declaración tributaria. *(Formal).*
\end{enumerate}
\begin{solucionBox}[Respuesta Oficial: Opción C]
\textbf{Justificación Oficial / Fundamentación:}\\[0.1cm]

\end{solucionBox}
\vspace{0.35cm}

\noindent\textbf{\sffamily Pregunta A.14:} El sustituto, cuando la ley lo prevea:\\[0.15cm]
\begin{enumerate}[label=\textbf{\alph*)},leftmargin=1.8em,itemsep=1.5pt,topsep=2pt,partopsep=0pt,parsep=0pt]
    \item Es el sujeto pasivo que se coloca junto al contribuyente. *(Junto al deudor se coloca el responsable).*
    \item Es el sujeto pasivo que se coloca en lugar del contribuyente. *(Art. 36.3 LGT).*
    \item Es el sujeto pasivo que se coloca junto y en lugar del contribuyente.
    \item Es el obligado tributario que se coloca en lugar del sujeto pasivo.
\end{enumerate}
\begin{solucionBox}[Respuesta Oficial: Opción B]
\textbf{Justificación Oficial / Fundamentación:}\\[0.1cm]

\end{solucionBox}
\vspace{0.35cm}

\noindent\textbf{\sffamily Pregunta A.15:} Cuando una persona fallece, como regla general:\\[0.15cm]
\begin{enumerate}[label=\textbf{\alph*)},leftmargin=1.8em,itemsep=1.5pt,topsep=2pt,partopsep=0pt,parsep=0pt]
    \item Se transmiten a sus herederos las deudas tributarias pendientes. *(Art. 39.1 LGT).*
    \item Se transmiten a sus herederos las deudas y las sanciones tributarias pendientes. *(Las sanciones nunca se transmiten).*
    \item Se extinguen las deudas tributarias en la mayoría de los casos.
    \item No se transmiten las deudas tributarias sino que se satisfacen exclusivamente con el patrimonio de la herencia.
\end{enumerate}
\begin{solucionBox}[Respuesta Oficial: Opción A]
\textbf{Justificación Oficial / Fundamentación:}\\[0.1cm]

\end{solucionBox}
\vspace{0.35cm}

\newpage
\subsection{Modelo B: Preguntas Específicas del Modelo B}
\begin{enunciadoBox}[Preguntas Diferenciales del Modelo B]
Preguntas incorporadas en la variante B del examen oficial que evalúan la deuda tributaria, tributos estatales vs autonómicos y locales, tasas y contribuciones especiales.
\end{enunciadoBox}
\vspace{0.4cm}

\noindent\textbf{\sffamily Pregunta B.2:} Pueden integrar la deuda tributaria:\\[0.15cm]
\begin{enumerate}[label=\textbf{\alph*)},leftmargin=1.8em,itemsep=1.5pt,topsep=2pt,partopsep=0pt,parsep=0pt]
    \item La cuota y las obligaciones tributarias accesorias.
    \item La obligación de realizar pagos a cuenta y las obligaciones tributarias accesorias.
    \item La cuota, los intereses de demora y determinados recargos.
    \item Todas las respuestas anteriores son correctas. *(Art. 58.1 y 58.2 LGT).*
\end{enumerate}
\begin{solucionBox}[Respuesta Oficial: Opción D]
\textbf{Justificación Oficial / Fundamentación:}\\[0.1cm]

\end{solucionBox}
\vspace{0.35cm}

\noindent\textbf{\sffamily Pregunta B.4:} La extinción de las deudas tributarias:\\[0.15cm]
\noindent\textbf{Opción Oficial:} \textbf{d)} Puede producirse por el pago y por otras formas de extinción. *(Art. 59.1 LGT: pago, prescripción, compensación, condonación...).*\\[0.15cm]
\begin{solucionBox}[Respuesta Oficial: Opción D]
\textbf{Justificación Oficial / Fundamentación:}\\[0.1cm]

\end{solucionBox}
\vspace{0.35cm}

\noindent\textbf{\sffamily Pregunta B.5:} Los impuestos indirectos:\\[0.15cm]
\noindent\textbf{Opción Oficial:} \textbf{a)} Como regla general, aunque con excepciones, permiten la repercusión de la cuota tributaria sobre un tercero Y  d) Son además impuestos objetivos. *(La plantilla da por válidas A y D).*\\[0.15cm]
\begin{solucionBox}[Respuesta Oficial: Opción A]
\textbf{Justificación Oficial / Fundamentación:}\\[0.1cm]

\end{solucionBox}
\vspace{0.35cm}

\noindent\textbf{\sffamily Pregunta B.6:} Son obligados tributarios:\\[0.15cm]
\begin{enumerate}[label=\textbf{\alph*)},leftmargin=1.8em,itemsep=1.5pt,topsep=2pt,partopsep=0pt,parsep=0pt]
    \item El contribuyente y el sustituto.
    \item El retenedor.
    \item Las respuestas a y b son correctas. *(Art. 35.2 LGT).*
\end{enumerate}
\begin{solucionBox}[Respuesta Oficial: Opción C]
\textbf{Justificación Oficial / Fundamentación:}\\[0.1cm]

\end{solucionBox}
\vspace{0.35cm}

\noindent\textbf{\sffamily Pregunta B.7:} El hecho imponible de una tasa puede ser:\\[0.15cm]
\noindent\textbf{Opción Oficial:} \textbf{c)} La prestación de servicios públicos. *(Art. 2.2.a LGT).*\\[0.15cm]
\begin{solucionBox}[Respuesta Oficial: Opción C]
\textbf{Justificación Oficial / Fundamentación:}\\[0.1cm]

\end{solucionBox}
\vspace{0.35cm}

\noindent\textbf{\sffamily Pregunta B.8:} Pertenecen al sistema impositivo estatal:\\[0.15cm]
\begin{enumerate}[label=\textbf{\alph*)},leftmargin=1.8em,itemsep=1.5pt,topsep=2pt,partopsep=0pt,parsep=0pt]
    \item El IVA y el IBI. *(El IBI es local).*
    \item El Impuesto de Transmisiones Patrimoniales y el Impuesto de Sucesiones y Donaciones, aunque estén cedidos a las Comunidades Autónomas. *(Aunque estén cedidos al 100\%, su titularidad pertenece al sistema impositivo estatal).*
    \item El IVA, el Impuesto sobre Construcciones, Instalaciones y Obras, y el Impuesto sobre Transmisiones Patrimoniales. *(El ICIO es local).*
\end{enumerate}
\begin{solucionBox}[Respuesta Oficial: Opción B]
\textbf{Justificación Oficial / Fundamentación:}\\[0.1cm]

\end{solucionBox}
\vspace{0.35cm}

\noindent\textbf{\sffamily Pregunta B.10:} Por el asfaltado de una calle y la consiguiente revalorización de los inmuebles situados en ella, se pagaría:\\[0.15cm]
\noindent\textbf{Opción Oficial:} \textbf{d)} Una contribución especial. *(Art. 2.2.b LGT).*\\[0.15cm]
\begin{solucionBox}[Respuesta Oficial: Opción D]
\textbf{Justificación Oficial / Fundamentación:}\\[0.1cm]

\end{solucionBox}
\vspace{0.35cm}

\noindent\textbf{\sffamily Pregunta B.13:} Son obligaciones tributarias materiales:\\[0.15cm]
\begin{enumerate}[label=\textbf{\alph*)},leftmargin=1.8em,itemsep=1.5pt,topsep=2pt,partopsep=0pt,parsep=0pt]
    \item Las sanciones tributarias.
    \item La obligación de llevar libros de contabilidad.
    \item El pago de la cuota. *(Art. 19 LGT: obligación tributaria principal).*
    \item La obligación de presentar la declaración tributaria.
\end{enumerate}
\begin{solucionBox}[Respuesta Oficial: Opción C]
\textbf{Justificación Oficial / Fundamentación:}\\[0.1cm]

\end{solucionBox}
\vspace{0.35cm}

\newpage
\section{Bloque II: Simulacro Oficial Parcial 9 Noviembre}
\label{sec:simulacro}

\begin{enunciadoBox}[Estructura y Reglas del Parcial Oficial de Aula]
\textbf{Convocatoria:} Parcial de Aula D16 (1,5 puntos sobre nota final, eliminatorio a partir de 0,75 puntos).\\
\textbf{Estructura:} 15 preguntas tipo test de 4 opciones (10 teóricas de LGT e IRPF $+$ 5 prácticas de liquidación).\\
\textbf{Criterio de corrección:} Cada 3 errores restan 1 acierto ($-0,33$ por fallo).
\end{enunciadoBox}

\vspace{0.3cm}
\subsection{Parte 1: 10 Preguntas Teóricas Clave}

\noindent\textbf{\sffamily Pregunta S.1:} [Teoría T1] Según el artículo 58.3 de la Ley General Tributaria, ¿cuál de los siguientes conceptos NO forma parte en ningún caso de la deuda tributaria?\\[0.15cm]
\begin{enumerate}[label=\textbf{\alph*)},leftmargin=1.8em,itemsep=1.5pt,topsep=2pt,partopsep=0pt,parsep=0pt]
    \item Los recargos por declaración extemporánea sin requerimiento previo (art. 27 LGT).
    \item Las sanciones tributarias pecuniarias.
    \item Los intereses de demora (art. 26 LGT).
    \item Los recargos del período ejecutivo (art. 28 LGT).
\end{enumerate}
\begin{solucionBox}[Respuesta Oficial: Opción B]
\textbf{Justificación Oficial / Fundamentación:}\\[0.1cm]
Art. 58.3 LGT: Las sanciones tributarias no formarán parte de la deuda tributaria en ningún caso. Sí forman parte la cuota líquida/diferencial y las obligaciones accesorias.
\end{solucionBox}
\vspace{0.35cm}

\noindent\textbf{\sffamily Pregunta S.2:} [Teoría T1] ¿Qué figura tributaria se exige al propietario de un inmueble cuando una obra pública municipal (ej. asfaltado y alumbrado nuevo) revaloriza especialmente su propiedad?\\[0.15cm]
\begin{enumerate}[label=\textbf{\alph*)},leftmargin=1.8em,itemsep=1.5pt,topsep=2pt,partopsep=0pt,parsep=0pt]
    \item Una tasa por aprovechamiento del dominio público.
    \item Un precio público municipal.
    \item Una contribución especial (art. 2.2.b LGT), con un importe máximo conjunto del 90\% del coste de la obra para el ayuntamiento.
    \item El Impuesto sobre Construcciones, Instalaciones y Obras (ICIO).
\end{enumerate}
\begin{solucionBox}[Respuesta Oficial: Opción C]
\textbf{Justificación Oficial / Fundamentación:}\\[0.1cm]
Art. 2.2.b LGT y art. 31 TRLRHL: Las contribuciones especiales gravan el beneficio o aumento de valor derivado de obras públicas o ampliación de servicios públicos, con un límite del 90\% del coste soportado.
\end{solucionBox}
\vspace{0.35cm}

\noindent\textbf{\sffamily Pregunta S.3:} [Teoría T1] En el sistema tributario local español (art. 59 TRLRHL), ¿cuáles son los impuestos obligatorios y cuáles los potestativos?\\[0.15cm]
\begin{enumerate}[label=\textbf{\alph*)},leftmargin=1.8em,itemsep=1.5pt,topsep=2pt,partopsep=0pt,parsep=0pt]
    \item Todos los impuestos municipales son de exigencia obligatoria.
    \item Son obligatorios el IBI, el IAE y el IVTM (todos directos); y potestativos el ICIO (indirecto) y el IIVTNU (directo).
    \item El IBI y el ICIO son obligatorios, y el IAE es voluntario.
    \item Solo el IBI es obligatorio en municipios menores de 20.000 habitantes.
\end{enumerate}
\begin{solucionBox}[Respuesta Oficial: Opción B]
\textbf{Justificación Oficial / Fundamentación:}\\[0.1cm]
Art. 59 TRLRHL: Los 3 impuestos de exigencia obligatoria en todos los consistorios son el IBI, el IAE y el IVTM. Los 2 potestativos (mediante ordenanza fiscal) son el ICIO y el IIVTNU.
\end{solucionBox}
\vspace{0.35cm}

\noindent\textbf{\sffamily Pregunta S.4:} [Teoría T1] ¿Qué diferencia jurídica fundamental existe entre el contribuyente y el sustituto del contribuyente (art. 36 LGT)?\\[0.15cm]
\begin{enumerate}[label=\textbf{\alph*)},leftmargin=1.8em,itemsep=1.5pt,topsep=2pt,partopsep=0pt,parsep=0pt]
    \item El contribuyente es siempre una persona jurídica y el sustituto una persona física.
    \item Ambos son sujetos pasivos, pero el contribuyente realiza el hecho imponible mientras que el sustituto queda obligado por ley en lugar de aquél a cumplir la obligación tributaria principal y formal.
    \item El sustituto solo responde de forma subsidiaria tras fallida del contribuyente.
    \item El sustituto es el recaudador público municipal.
\end{enumerate}
\begin{solucionBox}[Respuesta Oficial: Opción B]
\textbf{Justificación Oficial / Fundamentación:}\\[0.1cm]
Art. 36.1, 36.2 y 36.3 LGT: Sujeto pasivo engloba al contribuyente (quien realiza el hecho imponible) y al sustituto (quien por mandato legal se coloca en su lugar, con derecho a exigirle el reembolso).
\end{solucionBox}
\vspace{0.35cm}

\noindent\textbf{\sffamily Pregunta S.5:} [Teoría T1] En la liquidación de un tributo, ¿qué magnitud se obtiene al restar a la Cuota Íntegra las deducciones y bonificaciones legalmente aplicables (art. 56 LGT)?\\[0.15cm]
\begin{enumerate}[label=\textbf{\alph*)},leftmargin=1.8em,itemsep=1.5pt,topsep=2pt,partopsep=0pt,parsep=0pt]
    \item La Base Liquidable.
    \item La Cuota Diferencial.
    \item La Cuota Líquida.
    \item La Deuda Tributaria neta.
\end{enumerate}
\begin{solucionBox}[Respuesta Oficial: Opción C]
\textbf{Justificación Oficial / Fundamentación:}\\[0.1cm]
Art. 56 LGT: Cuota Íntegra - Deducciones/Bonificaciones = Cuota Líquida. Cuota Líquida - Pagos a cuenta = Cuota Diferencial.
\end{solucionBox}
\vspace{0.35cm}

\noindent\textbf{\sffamily Pregunta S.6:} [Teoría IRPF] Si un contribuyente traslada su residencia fiscal a un territorio calificado como jurisdicción no cooperativa (paraíso fiscal), según el artículo 8.2 de la LIRPF:\\[0.15cm]
\begin{enumerate}[label=\textbf{\alph*)},leftmargin=1.8em,itemsep=1.5pt,topsep=2pt,partopsep=0pt,parsep=0pt]
    \item Deja de tributar en España de forma inmediata e irrevocable.
    \item Seguirá teniendo la condición de contribuyente del IRPF en el período impositivo en que efectúe el cambio y durante los cuatro períodos impositivos siguientes.
    \item Tributa únicamente por obligación real en el IRNR.
    \item Paga una tasa fija única del 35\% sobre su renta mundial.
\end{enumerate}
\begin{solucionBox}[Respuesta Oficial: Opción B]
\textbf{Justificación Oficial / Fundamentación:}\\[0.1cm]
Art. 8.2 LIRPF: La regla especial de cuarentena fiscal antiparaísos mantiene la condición de contribuyente por obligación personal (IRPF) en el año del traslado y los 4 siguientes (5 años en total).
\end{solucionBox}
\vspace{0.35cm}

\noindent\textbf{\sffamily Pregunta S.7:} [Teoría IRPF] En el cómputo de los 183 días de permanencia para determinar la residencia habitual en España (art. 9.1.a LIRPF), las ausencias esporádicas:\\[0.15cm]
\begin{enumerate}[label=\textbf{\alph*)},leftmargin=1.8em,itemsep=1.5pt,topsep=2pt,partopsep=0pt,parsep=0pt]
    \item Se descuentan en todo caso de los días de estancia en España.
    \item Se suman a los días de permanencia en España, salvo que el contribuyente acredite su residencia fiscal en otro país mediante certificado de la autoridad fiscal correspondiente.
    \item Solo se computan si son viajes de vacaciones de menos de 10 días.
    \item Quedan excluidas si el contribuyente aporta billetes de avión y visados turísticos.
\end{enumerate}
\begin{solucionBox}[Respuesta Oficial: Opción B]
\textbf{Justificación Oficial / Fundamentación:}\\[0.1cm]
Art. 9.1.a LIRPF: Para determinar el período de permanencia se computarán las ausencias esporádicas, salvo que el contribuyente acredite su residencia fiscal en otro país mediante certificado oficial.
\end{solucionBox}
\vspace{0.35cm}

\noindent\textbf{\sffamily Pregunta S.8:} [Teoría IRPF] ¿Cuál es el límite legal máximo anual de deducción por cuotas satisfechas a colegios profesionales cuando la colegiación sea obligatoria (art. 19.2.d LIRPF)?\\[0.15cm]
\begin{enumerate}[label=\textbf{\alph*)},leftmargin=1.8em,itemsep=1.5pt,topsep=2pt,partopsep=0pt,parsep=0pt]
    \item 300 \euro{} anuales.
    \item 500 \euro{} anuales.
    \item 1.000 \euro{} anuales.
    \item El 100\% de la cuota satisfecha sin límite.
\end{enumerate}
\begin{solucionBox}[Respuesta Oficial: Opción B]
\textbf{Justificación Oficial / Fundamentación:}\\[0.1cm]
Art. 19.2.d LIRPF y art. 10 RIRPF: Las cuotas de colegios profesionales obligatorios para el ejercicio del puesto son deducibles con un límite máximo de 500 \euro{} anuales.
\end{solucionBox}
\vspace{0.35cm}

\noindent\textbf{\sffamily Pregunta S.9:} [Teoría IRPF] En el arrendamiento de bienes inmuebles destinados a vivienda, ¿qué límite conjunto tienen los gastos de intereses de hipoteca y de conservación/reparación (art. 23.1.a LIRPF)?\\[0.15cm]
\begin{enumerate}[label=\textbf{\alph*)},leftmargin=1.8em,itemsep=1.5pt,topsep=2pt,partopsep=0pt,parsep=0pt]
    \item No tienen ningún límite, pudiendo arrojar un rendimiento negativo ilimitado.
    \item No pueden superar la cuantía de los rendimientos íntegros obtenidos de ese inmueble; el exceso no deducido en el año puede compensarse en los cuatro años siguientes.
    \item Tienen como límite el 50\% de los ingresos de alquiler.
    \item Solo se permite deducir hasta 1.000 \euro{} al año por contrato.
\end{enumerate}
\begin{solucionBox}[Respuesta Oficial: Opción B]
\textbf{Justificación Oficial / Fundamentación:}\\[0.1cm]
Art. 23.1.a LIRPF: Los gastos de financiación y de conservación/reparación tienen como tope los rendimientos íntegros del propio inmueble. El exceso se deduce en los 4 ejercicios siguientes.
\end{solucionBox}
\vspace{0.35cm}

\noindent\textbf{\sffamily Pregunta S.10:} [Teoría IRPF] ¿Qué porcentaje de imputación de rentas inmobiliarias se aplica a una segunda residencia urbana a disposición del propietario cuyo valor catastral fue revisado en los últimos 10 años (art. 85 LIRPF)?\\[0.15cm]
\begin{enumerate}[label=\textbf{\alph*)},leftmargin=1.8em,itemsep=1.5pt,topsep=2pt,partopsep=0pt,parsep=0pt]
    \item 2\% sobre el valor catastral.
    \item 1,1\% sobre el valor catastral.
    \item 1,1\% sobre el 50\% del valor de adquisición.
    \item 0,5\% sobre el valor de mercado.
\end{enumerate}
\begin{solucionBox}[Respuesta Oficial: Opción B]
\textbf{Justificación Oficial / Fundamentación:}\\[0.1cm]
Art. 85.1 LIRPF: El porcentaje general es del 2\%, pero se reduce al 1,1\% en caso de inmuebles cuyos valores catastrales hayan sido revisados en los últimos 10 períodos impositivos.
\end{solucionBox}
\vspace{0.35cm}

\newpage
\subsection{Parte 2: Caso Práctico de Liquidación --- Dña. Elena}

\begin{enunciadoBox}[Enunciado Oficial del Caso Práctico]
Dña. Elena (42 años), residente en Granada, presenta los siguientes datos fiscales en el ejercicio:
\begin{enumerate}[label=\textbf{\arabic*)}]
    \item Trabaja como directiva en una empresa desde hace 6 años con un sueldo bruto anual de $44.000\text{ \euro}$, habiendo cotizado a la Seguridad Social $2.650\text{ \euro}$ y pagado $580\text{ \euro}$ al Colegio Oficial de Economistas (colegiación obligatoria para su puesto) y $200\text{ \euro}$ a un sindicato.
    \item En mayo se desplazó en su vehículo particular de Granada a Sevilla ($500\text{ km}$ ida y vuelta en total) durante 3 días (pernoctando 2 noches en hotel y regresando el 3.er día). La empresa le abonó: $0,36\text{ \euro/km}$ ($180\text{ \euro}$), peajes y parking justificados ($40\text{ \euro}$), factura de hotel de las 2 noches ($210\text{ \euro}$) y $220\text{ \euro}$ en concepto de dieta de manutención por los 3 días.
    \item Además, la empresa le cedió para uso exclusivamente particular un vehículo propiedad de la empresa cuyo coste de adquisición fue de $28.000\text{ \euro}$ (sin distintivo de eficiencia energética). El ingreso a cuenta ($20\%$) fue REPERCUTIDO a la trabajadora.
    \item Es propietaria de un piso en Granada (Coste de adquisición satisfecho $= 180.000\text{ \euro}$; Valor Catastral $= 90.000\text{ \euro}$, correspondiendo el $20\%$ al suelo y el $80\%$ a la construcción), alquilado todo el año como vivienda habitual a un matrimonio por $750\text{ \euro/mes}$ ($9.000\text{ \euro/año}$) mediante contrato firmado en 2024 (reducción general del $50\%$). Gastos anuales soportados: intereses de hipoteca $5.500\text{ \euro}$, reparación de caldera y pintura $4.000\text{ \euro}$, IBI $450\text{ \euro}$ y comunidad $650\text{ \euro}$.
    \item Posee un apartamento de vacaciones en Salobreña (Valor Catastral revisado hace 4 años $= 80.000\text{ \euro}$) que ha permanecido a su disposición todo el año (vacío).
\end{enumerate}
\end{enunciadoBox}

\vspace{0.3cm}

\noindent\textbf{\sffamily Pregunta S.11:} [Caso Práctico Dña. Elena — Apartado 2] Por el desplazamiento en su vehículo de Granada a Sevilla (500 km, 3 días, 2 noches en hotel), la empresa abonó 180 \euro{} de kilometraje (0,36 \euro{}/km), 40 \euro{} de peaje/parking, 210 \euro{} de hotel y 220 \euro{} de manutención. ¿Cuánto tributa Dña. Elena por dietas?\\[0.15cm]
\begin{enumerate}[label=\textbf{\alph*)},leftmargin=1.8em,itemsep=1.5pt,topsep=2pt,partopsep=0pt,parsep=0pt]
    \item 0 \euro{}, porque todo el viaje obedece a razones laborales justificadas.
    \item 136,65 \euro{} como rendimiento dinerario del trabajo (50 \euro{} por exceso de kilometraje + 86,65 \euro{} por exceso de manutención).
    \item 220 \euro{} íntegros por la dieta de manutención.
    \item 50 \euro{} exclusivamente por el exceso de kilometraje.
\end{enumerate}
\begin{solucionBox}[Respuesta Oficial: Opción B]
\textbf{Justificación Oficial / Fundamentación:}\\[0.1cm]
Art. 9 RIRPF: 1) Locomoción: 500 km $\times$  0,26 \euro{}/km = 130 \euro{} exentos. Cobró 180 \euro{} -> exceso sujeto = 50 \euro{}. Peaje y parking (40 \euro{}) y hotel (210 \euro{}) justificados -> 100\% exentos. 2) Manutención: 2 noches hotel = 2 días con pernocta a 53,34 \euro{} (106,68 \euro{}) + 1 día de vuelta sin pernocta a 26,67 \euro{} = 133,35 \euro{} exentos. Cobró 220 \euro{} -> exceso sujeto = 220 \euro{} - 133,35 \euro{} = 86,65 \euro{}. Total sujeto = 50 \euro{} + 86,65 \euro{} = 136,65 \euro{}.
\end{solucionBox}
\vspace{0.35cm}

\noindent\textbf{\sffamily Pregunta S.12:} [Caso Práctico Dña. Elena — Apartado 3] Por la cesión para uso exclusivamente particular de un vehículo propiedad de la empresa (coste de compra 28.000 \euro{} sin distintivo ambiental) con ingreso a cuenta del 20\% repercutido a la trabajadora, ¿cuál es el rendimiento íntegro del trabajo en especie?\\[0.15cm]
\begin{enumerate}[label=\textbf{\alph*)},leftmargin=1.8em,itemsep=1.5pt,topsep=2pt,partopsep=0pt,parsep=0pt]
    \item 5.600 \euro{} anuales (20\% del coste de adquisición), sin añadir el ingreso a cuenta al haber sido repercutido.
    \item 6.720 \euro{} anuales (5.600 \euro{} de valoración más 1.120 \euro{} de ingreso a cuenta).
    \item 2.800 \euro{} anuales (10\% del valor de compra).
    \item 0 \euro{}, porque el uso de vehículos de empresa está exento.
\end{enumerate}
\begin{solucionBox}[Respuesta Oficial: Opción A]
\textbf{Justificación Oficial / Fundamentación:}\\[0.1cm]
Art. 43.1.1ª.b y 43.2 LIRPF: La valoración en especie del uso de un vehículo propiedad de la empresa es del 20\% anual del coste de adquisición satisfecho (20\% $\times$  28.000 \euro{} = 5.600 \euro{}). Como el ingreso a cuenta fue repercutido a la trabajadora descontándoselo, no se adiciona a la valoración: rendimiento íntegro = 5.600 \euro{}.
\end{solucionBox}
\vspace{0.35cm}

\noindent\textbf{\sffamily Pregunta S.13:} [Caso Práctico Dña. Elena — Apartados 1, 2 y 3] ¿Cuál es el Rendimiento Neto del Trabajo Personal (RNTP) de Dña. Elena antes de reducciones del artículo 20?\\[0.15cm]
\begin{enumerate}[label=\textbf{\alph*)},leftmargin=1.8em,itemsep=1.5pt,topsep=2pt,partopsep=0pt,parsep=0pt]
    \item 44.386,65 \euro{}.
    \item 45.000,00 \euro{}.
    \item 42.250,00 \euro{}.
    \item 47.186,65 \euro{}.
\end{enumerate}
\begin{solucionBox}[Respuesta Oficial: Opción A]
\textbf{Justificación Oficial / Fundamentación:}\\[0.1cm]
Arts. 17 y 19 LIRPF: Rendimientos íntegros = Sueldo dinerario (44.000 \euro{}) + Exceso dietas sujetas (136,65 \euro{}) + Retribución especie vehículo (5.600 \euro{}) = 49.736,65 \euro{}. Gastos deducibles: Seguridad Social (2.650 \euro{}) + Colegio Oficial obligatorio con tope legal (500 \euro{}, art. 19.2.d) + Cuota sindicato (200 \euro{}, art. 19.2.d) + Otros gastos generales (2.000 \euro{}, art. 19.2.f) = 5.350 \euro{}. RNTP previo = 49.736,65 \euro{} - 5.350 \euro{} = 44.386,65 \euro{}.
\end{solucionBox}
\vspace{0.35cm}

\noindent\textbf{\sffamily Pregunta S.14:} [Caso Práctico Dña. Elena — Apartado 4] Por el piso alquilado en Granada como vivienda habitual (contrato de 2024 a 750 \euro{}/mes = 9.000 \euro{}/año; intereses 5.500 \euro{}, reparación 4.000 \euro{}, IBI 450 \euro{}, comunidad 650 \euro{}, amortización 3\% s/ construcción de 144.000 \euro{} = 4.320 \euro{}), ¿cuál es el Rendimiento Neto Reducido del Capital Inmobiliario?\\[0.15cm]
\begin{enumerate}[label=\textbf{\alph*)},leftmargin=1.8em,itemsep=1.5pt,topsep=2pt,partopsep=0pt,parsep=0pt]
    \item -5.420 \euro{} (al ser rendimiento neto negativo, no procede aplicar la reducción del 50\%).
    \item -2.710 \euro{} (aplicando la reducción del 50\% por vivienda habitual).
    \item 0 \euro{} (los gastos deducibles nunca pueden dar lugar a rendimiento negativo).
    \item -420 \euro{} (computando todos los intereses y reparaciones sin límite).
\end{enumerate}
\begin{solucionBox}[Respuesta Oficial: Opción A]
\textbf{Justificación Oficial / Fundamentación:}\\[0.1cm]
Arts. 22 y 23 LIRPF: 1) Ingresos íntegros = 9.000 \euro{}. 2) Gastos de reparación y financiación con tope: 5.500 \euro{} + 4.000 \euro{} = 9.500 \euro{}. Se deducen hasta el límite de los ingresos íntegros (9.000 \euro{}); el exceso de 500 \euro{} queda pendiente para los 4 años siguientes. Saldo tras gastos con límite = 0 \euro{}. 3) Gastos necesarios sin tope: IBI (450 \euro{}) + Comunidad (650 \euro{}) + Amortización (3\% $\times$  180.000 $\times$  80\% = 4.320 \euro{}) = 5.420 \euro{}. Rendimiento neto = 0 \euro{} - 5.420 \euro{} = -5.420 \euro{}. 4) Reducción por arrendamiento de vivienda (art. 23.2 LIRPF): Al tratarse de un rendimiento neto negativo, no es aplicable ninguna reducción (el art. 23.2 se reserva expresamente a rendimientos netos positivos). Por tanto, el Rendimiento Neto Reducido es de -5.420 \euro{}.
\end{solucionBox}
\vspace{0.35cm}

\noindent\textbf{\sffamily Pregunta S.15:} [Caso Práctico Dña. Elena — Apartado 5] ¿Qué renta debe declarar Dña. Elena por el apartamento de vacaciones en Salobreña (Valor Catastral revisado hace 4 años = 80.000 \euro{}) que estuvo vacío todo el año?\\[0.15cm]
\begin{enumerate}[label=\textbf{\alph*)},leftmargin=1.8em,itemsep=1.5pt,topsep=2pt,partopsep=0pt,parsep=0pt]
    \item 1.600 \euro{} (2\% del valor catastral).
    \item 880 \euro{} como Imputación de Rentas Inmobiliarias (1,1\% de 80.000 \euro{} al estar revisado en los últimos 10 períodos impositivos, art. 85 LIRPF) sin poder deducir ningún gasto.
    \item 0 \euro{}, porque al no haber generado ingresos no tributa.
    \item 880 \euro{} menos el IBI satisfecho.
\end{enumerate}
\begin{solucionBox}[Respuesta Oficial: Opción B]
\textbf{Justificación Oficial / Fundamentación:}\\[0.1cm]
Art. 85 LIRPF: Inmuebles urbanos no afectos ni alquilados generan imputación de rentas al 1,1\% sobre el valor catastral revisado: 1,1\% $\times$  80.000 \euro{} = 880 \euro{}. No se admite la deducción de ningún gasto.
\end{solucionBox}
\vspace{0.35cm}

\newpage
\section{Bloque III: Examen Oficial Convocatoria Ordinaria 2019}
\label{sec:examen2019}

\begin{enunciadoBox}[Enunciado Oficial de la Profesora: Matrimonio Casares --- Gutiérrez]
Doña María Casares y D. Julio Gutiérrez, ambos de 50 años, son un matrimonio casado en régimen de gananciales desde el año 2000 y residentes en Granada. Tienen dos hijos solteros de 19 y 17 años, que están estudiando, sin rentas y que conviven con ellos. También reside con el matrimonio la madre de Doña María de 76 años, sin rentas y con un grado de minusvalía reconocido del $40\%$.

\textbf{Actividad económica de Doña María:} Tiene una escuela de patronaje y corte y confección (IAE 933.9) que montó en 2016, por la que está en el régimen de módulos en el IRPF. No ha contratado a nadie hasta 2018 que ha tenido una persona asalariada a jornada completa desde el 1 de Febrero (1650 horas). La actividad se desarrolla en un local de $40\text{ m}^2$, adquirido en 2017 por $90.000\text{ \euro}$ (de los que $45.000\text{ \euro}$ corresponden al suelo) y se sabe que, en 2017, se compraron máquinas de coser por $2.000\text{ \euro}$ (valor por máquina de $300\text{ \euro}$) y un ordenador de $2.000\text{ \euro}$. En el ejercicio 2018, se han comprado más máquinas de coser por $1.000\text{ \euro}$, por $260\text{ \euro}$ cada una, y dos máquinas remalladoras profesionales de segunda mano, cada una de las cuales ha supuesto una inversión de $700\text{ \euro}$. En el 2017, se obtuvieron ingresos por $20.000\text{ \euro}$ y en 2018, por $50.000\text{ \euro}$. Se viene amortizando a coeficientes máximos y los pagos fraccionados han sido de $456\text{ \euro}$.

\textbf{Rendimientos del trabajo de D. Julio:} D. Julio es ingeniero y trabaja para una empresa tecnológica. Su salario íntegro en 2018 ha sido de $62.000\text{ \euro}$ (tipo de retención $25\%$, retenciones $15.500\text{ \euro}$, y Seguridad Social $2.900\text{ \euro}$). En mayo de 2018 asistió a un Congreso de Barcelona, durante 3 días (dos noches) del que D. Julio conserva factura del hotel donde se quedó por importe de $300\text{ \euro}$ y un billete de avión por $90\text{ \euro}$. Por la asistencia al congreso, la empresa le paga $400\text{ \euro}$.
Se sabe que en 2018, D. Julio recibió un premio por sus 10 años al servicio de la empresa por $5.000\text{ \euro}$ (de los que $875\text{ \euro}$ corresponden a retención), que la empresa le paga una prima anual por un seguro de enfermedad de $600\text{ \euro}$ y que se han realizado aportaciones a un plan de pensiones por $6.000\text{ \euro}$, de los que $3.000\text{ \euro}$ corresponden a la empresa. Los derechos consolidados en el Plan de Pensiones ascienden a $40.000\text{ \euro}$.

\textbf{Patrimonio e información patrimonial adicional en 2018:}
\begin{itemize}[itemsep=1.5pt]
    \item Residen en una vivienda adquirida en 2001 por $180.000\text{ \euro}$, cuyo valor catastral es de $200.000\text{ \euro}$ ($50\%$ suelo). IBI: $312\text{ \euro}$.
    \item El 1 de abril de 2018, Doña María vendió una finca rústica, heredada de su padre en 2013, por $30.000\text{ \euro}$. Valor en Sucesiones: $35.000\text{ \euro}$, valor catastral revisado $23.000\text{ \euro}$ ($20\%$ suelo). La finca ha estado improductiva todo este tiempo.
    \item En 2010 adquirieron un apartamento en la playa, alquilado en 2018 todo el año por $400\text{ \euro/mes}$. Valor adquisición: $130.000\text{ \euro}$, VC revisado $90.000\text{ \euro}$ ($25\%$ suelo). IBI asumido: $130\text{ \euro}$.
    \item En enero de 2018, por fallecimiento de un primo lejano de sus padres, D. Julio heredó 300 acciones cotizadas en bolsa (valor en ISD $= 21.000\text{ \euro}$, es decir $70\text{ \euro/acción}$). Dividendos líquidos: $2.430\text{ \euro}$ (retención $570\text{ \euro}$), gastos administración $5\text{ \euro}$.
    \item De las 300 acciones, vendió 50 en noviembre a $100\text{ \euro/acción}$. Cotización a 31 de diciembre: $101\text{ \euro}$.
    \item Cuenta corriente conjunta con saldo medio último trimestre de $62.000\text{ \euro}$ y saldo a 31 dic de $75.000\text{ \euro}$.
    \item Doña María tiene bonos a largo plazo heredados de su padre con intereses de $500\text{ \euro}$ en 2018 (cotización último trimestre: $12.000\text{ \euro}$).
    \item Donativos de D. Julio a Cruz Roja: $300\text{ \euro}$.
\end{itemize}
\end{enunciadoBox}

\vspace{0.4cm}

\noindent\textbf{\sffamily Pregunta 2019.1:} Las acciones heredadas por D. Julio de su tío:\\[0.15cm]
\begin{enumerate}[label=\textbf{\alph*)},leftmargin=1.8em,itemsep=1.5pt,topsep=2pt,partopsep=0pt,parsep=0pt]
    \item Suponen una ganancia de capital en el IRPF de D. Julio y un incremento de patrimonio a efectos del Impuesto sobre Sucesiones y Donaciones.
    \item No están sujetas al IRPF de D. Julio por su adquisición.
    \item No generan rentas a efectos del IRPF del ejercicio 2018.
    \item No supone una adquisición sujeta al ISD, por tratarse de herencia familiar.
\end{enumerate}
\begin{solucionBox}[Respuesta Oficial: Opción B]
\textbf{Justificación Oficial / Fundamentación:}\\[0.1cm]
Art. 6.4 LIRPF: La adquisición gratuita de bienes por herencia (mortis causa) no está sujeta al IRPF, tributando exclusivamente en el Impuesto sobre Sucesiones y Donaciones (ISD).
\end{solucionBox}
\vspace{0.35cm}

\noindent\textbf{\sffamily Pregunta 2019.2:} Tendría el matrimonio que presentar declaración por el Impuesto sobre el Patrimonio:\\[0.15cm]
\begin{enumerate}[label=\textbf{\alph*)},leftmargin=1.8em,itemsep=1.5pt,topsep=2pt,partopsep=0pt,parsep=0pt]
    \item El matrimonio no, pero D. Julio sí.
    \item El matrimonio no, pero Doña María si.
    \item Si por que el importe de su patrimonio neto es superior a 600.000\euro{}.
    \item Ninguna es correcta.
\end{enumerate}
\begin{solucionBox}[Respuesta Oficial: Opción D]
\textbf{Justificación Oficial / Fundamentación:}\\[0.1cm]
Art. 37 Ley del Impuesto sobre el Patrimonio: El mínimo exento en Andalucía es de 700.000 \euro{} por contribuyente más 300.000 \euro{} de exención por vivienda habitual, y la obligación formal de declarar sin cuota a ingresar exige superar 2.000.000 \euro{} de patrimonio bruto. Como ninguno de los dos supera estos límites, ninguno está obligado a declarar (por tanto, 'Ninguna es correcta').
\end{solucionBox}
\vspace{0.35cm}

\noindent\textbf{\sffamily Pregunta 2019.3:} Cuál sería la reducción por incentivos al empleo en la actividad de Dña María:\\[0.15cm]
\begin{enumerate}[label=\textbf{\alph*)},leftmargin=1.8em,itemsep=1.5pt,topsep=2pt,partopsep=0pt,parsep=0pt]
    \item 0,364
    \item 456,57
    \item 0,091
    \item 114,06
\end{enumerate}
\begin{solucionBox}[Respuesta Oficial: Opción A]
\textbf{Justificación Oficial / Fundamentación:}\\[0.1cm]
Orden de Módulos (IRPF): En estimación objetiva, el coeficiente por incentivos al empleo se calcula en función del incremento de plantilla de personal asalariado (1 empleado a jornada completa desde el 1 de febrero = 11/12 = 0,9167 personas/año, aplicando los tramos del módulo de empleo resulta el coeficiente 0,364).
\end{solucionBox}
\vspace{0.35cm}

\noindent\textbf{\sffamily Pregunta 2019.4:} Cuál sería el rendimiento neto del trabajo personal de D. Julio:\\[0.15cm]
\begin{enumerate}[label=\textbf{\alph*)},leftmargin=1.8em,itemsep=1.5pt,topsep=2pt,partopsep=0pt,parsep=0pt]
    \item 65.225
    \item 58.225
    \item 63.725
    \item 63.700
\end{enumerate}
\begin{solucionBox}[Respuesta Oficial: Opción D]
\textbf{Justificación Oficial / Fundamentación:}\\[0.1cm]
Arts. 17 a 19 LIRPF: Salario íntegro (62.000 \euro{}) + Premio antigüedad 5.000 \euro{} con reducción 30\% irregularidad = 3.500 \euro{} + Seguro enfermedad (600 - 500 exento = 100 \euro{}) + Plan pensiones empresa (3.000 \euro{}) + Dietas sujetas por exceso (10 \euro{} según pregunta 17) = 68.610 \euro{}. Gastos deducibles: Seguridad Social (2.900 \euro{}) + Otros gastos generales (2.000 \euro{}) = 4.900 \euro{}. Rendimiento neto = 68.610 \euro{} - 4.900 \euro{} = 63.710 \euro{} $\approx$  63.700 \euro{}.
\end{solucionBox}
\vspace{0.35cm}

\noindent\textbf{\sffamily Pregunta 2019.5:} Cual seria el rendimiento neto de módulos de la actividad de Doña Maria:\\[0.15cm]
\begin{enumerate}[label=\textbf{\alph*)},leftmargin=1.8em,itemsep=1.5pt,topsep=2pt,partopsep=0pt,parsep=0pt]
    \item 12.103,74
    \item 12.740,77
    \item 12.402,99
    \item 13.448,6
\end{enumerate}
\begin{solucionBox}[Respuesta Oficial: Opción A]
\textbf{Justificación Oficial / Fundamentación:}\\[0.1cm]
Orden de Módulos IAE 933.9: Rendimiento neto previo calculado con los módulos de personal no asalariado (titular), personal asalariado y superficie de local, minorado por los incentivos al empleo e inversión y amortizaciones a coeficientes máximos resulta 12.103,74 \euro{}.
\end{solucionBox}
\vspace{0.35cm}

\noindent\textbf{\sffamily Pregunta 2019.6:} El mínimo personal y familiar de Doña María es:\\[0.15cm]
\begin{enumerate}[label=\textbf{\alph*)},leftmargin=1.8em,itemsep=1.5pt,topsep=2pt,partopsep=0pt,parsep=0pt]
    \item 16.200\euro{}
    \item 12.500\euro{}
    \item 12.300\euro{}
    \item 13.650\euro{}
\end{enumerate}
\begin{solucionBox}[Respuesta Oficial: Opción C]
\textbf{Justificación Oficial / Fundamentación:}\\[0.1cm]
Arts. 57 a 60 LIRPF: Mínimo contribuyente (5.550 \euro{}) + 1er hijo al 50\% (1.200 \euro{}) + 2º hijo al 50\% (1.350 \euro{}) + Ascendiente > 65 años (1.150 \euro{}) + Discapacidad ascendiente 40\% (3.000 \euro{}) = 5.550 + 1.200 + 1.350 + 4.150 = 12.250 \euro{} $\approx$  12.300 \euro{}.
\end{solucionBox}
\vspace{0.35cm}

\noindent\textbf{\sffamily Pregunta 2019.7:} La finca rústica:\\[0.15cm]
\begin{enumerate}[label=\textbf{\alph*)},leftmargin=1.8em,itemsep=1.5pt,topsep=2pt,partopsep=0pt,parsep=0pt]
    \item No genera rentas hasta su venta
    \item Genera una imputación de rentas inmobiliarias por 253\euro{}
    \item Genera una imputación de rentas inmobiliarias por 63,25\euro{}
    \item Tributa en el impuesto de Sucesiones y Donaciones por su transmisión.
\end{enumerate}
\begin{solucionBox}[Respuesta Oficial: Opción A]
\textbf{Justificación Oficial / Fundamentación:}\\[0.1cm]
Art. 85.1 LIRPF: El régimen de imputación de rentas inmobiliarias grava exclusivamente bienes inmuebles URBANOS no arrendados ni afectos. Las fincas rústicas no generan imputación de rentas.
\end{solucionBox}
\vspace{0.35cm}

\noindent\textbf{\sffamily Pregunta 2019.8:} Los rendimientos netos del capital mobiliario de D. Julio han sido:\\[0.15cm]
\begin{enumerate}[label=\textbf{\alph*)},leftmargin=1.8em,itemsep=1.5pt,topsep=2pt,partopsep=0pt,parsep=0pt]
    \item 4.500\euro{}
    \item 4.495\euro{}
    \item 2.995\euro{}
    \item 2.425\euro{}
\end{enumerate}
\begin{solucionBox}[Respuesta Oficial: Opción C]
\textbf{Justificación Oficial / Fundamentación:}\\[0.1cm]
Arts. 25.1.a y 26.1.a LIRPF: Dividendos íntegros de las acciones privativas heredadas = 2.430 \euro{} netos + 570 \euro{} retención = 3.000 \euro{}. Gasto de custodia deducible = 5 \euro{}. Rendimiento neto = 3.000 \euro{} - 5 \euro{} = 2.995 \euro{}.
\end{solucionBox}
\vspace{0.35cm}

\noindent\textbf{\sffamily Pregunta 2019.9:} La base imponible del ahorro de Dña. Maria es:\\[0.15cm]
\begin{enumerate}[label=\textbf{\alph*)},leftmargin=1.8em,itemsep=1.5pt,topsep=2pt,partopsep=0pt,parsep=0pt]
    \item 0\euro{}
    \item -4.500\euro{}
    \item 125\euro{}
    \item 375\euro{}
\end{enumerate}
\begin{solucionBox}[Respuesta Oficial: Opción D]
\textbf{Justificación Oficial / Fundamentación:}\\[0.1cm]
Art. 49 LIRPF: Bonos heredados reportan 500 \euro{} de intereses (RCM). La venta de la finca rústica arroja una pérdida patrimonial de 30.000 - 35.000 = -5.000 \euro{}. En la base del ahorro, el saldo negativo de ganancias se compensa con el saldo positivo de RCM con el límite del 25\% (25\% de 500 = 125 \euro{}). Saldo neto positivo del ahorro = 500 - 125 = 375 \euro{}.
\end{solucionBox}
\vspace{0.35cm}

\noindent\textbf{\sffamily Pregunta 2019.10:} Por la venta de las acciones se obtiene:\\[0.15cm]
\begin{enumerate}[label=\textbf{\alph*)},leftmargin=1.8em,itemsep=1.5pt,topsep=2pt,partopsep=0pt,parsep=0pt]
    \item Un rendimiento neto del capital mobiliario por 1.500\euro{} para D. Julio.
    \item Una ganancia de capital de 1.500\euro{} para D. Julio.
    \item Un rendimiento neto del capital mobiliario para el matrimonio de 1.500\euro{}
    \item Una ganancia de capital para el matrimonio de 1.500\euro{}
\end{enumerate}
\begin{solucionBox}[Respuesta Oficial: Opción B]
\textbf{Justificación Oficial / Fundamentación:}\\[0.1cm]
Arts. 33 a 35 LIRPF: Venta de 50 acciones a 100 \euro{} (5.000 \euro{}) adquiridas en la herencia a 70 \euro{} (3.500 \euro{}). Ganancia patrimonial = 5.000 \euro{} - 3.500 \euro{} = 1.500 \euro{}. Al ser títulos privativos heredados por D. Julio (art. 11.5 LIRPF), la ganancia corresponde íntegramente a D. Julio.
\end{solucionBox}
\vspace{0.35cm}

\noindent\textbf{\sffamily Pregunta 2019.11:} El apartamento de la playa genera:\\[0.15cm]
\begin{enumerate}[label=\textbf{\alph*)},leftmargin=1.8em,itemsep=1.5pt,topsep=2pt,partopsep=0pt,parsep=0pt]
    \item Unos rendimientos netos de capital inmobiliario de 1.772,5\euro{} para cada uno de los cónyuges.
    \item Unos rendimientos netos de capital inmobiliario de 709\euro{} para cada uno de los cónyuges.
    \item Unos rendimientos netos de capital inmobiliario de 349\euro{} para cada uno de los cónyuges.
    \item Unos rendimientos neto del capital inmobiliario de 872,5\euro{} para cada uno de los cónyuges.
\end{enumerate}
\begin{solucionBox}[Respuesta Oficial: Opción D]
\textbf{Justificación Oficial / Fundamentación:}\\[0.1cm]
Arts. 22 y 23 LIRPF: Ingresos íntegros = 400 \euro{} $\times$  12 = 4.800 \euro{}. Gastos deducibles: IBI = 130 \euro{}. Amortización 3\% sobre el mayor valor (coste construcción = 130.000 \euro{} $\times$  75\% = 97.500 \euro{} vs VC construcción = 90.000 \euro{} $\times$  75\% = 67.500 \euro{}): 3\% de 97.500 \euro{} = 2.925 \euro{}. Rendimiento neto total = 4.800 \euro{} - 130 \euro{} - 2.925 \euro{} = 1.745 \euro{}. Al 50\% ganancial = 872,50 \euro{} para cada cónyuge.
\end{solucionBox}
\vspace{0.35cm}

\noindent\textbf{\sffamily Pregunta 2019.12:} La base liquidable general en tributación conjunta sería\\[0.15cm]
\begin{enumerate}[label=\textbf{\alph*)},leftmargin=1.8em,itemsep=1.5pt,topsep=2pt,partopsep=0pt,parsep=0pt]
    \item 65.916,37\euro{}
    \item 69.316,37\euro{}
    \item 68.916,37\euro{}
    \item 75.316,37\euro{}
\end{enumerate}
\begin{solucionBox}[Respuesta Oficial: Opción B]
\textbf{Justificación Oficial / Fundamentación:}\\[0.1cm]
Arts. 50 y 84 LIRPF: Integración de rendimientos del trabajo de D. Julio + rendimiento de módulos de Doña María + rendimientos de capital inmobiliario, minorados por la reducción por tributación conjunta (3.400 \euro{}) y la reducción por planes de pensiones (6.000 \euro{}) arroja 69.316,37 \euro{}.
\end{solucionBox}
\vspace{0.35cm}

\noindent\textbf{\sffamily Pregunta 2019.13:} La cuota íntegra estatal de Doña María es:\\[0.15cm]
\begin{enumerate}[label=\textbf{\alph*)},leftmargin=1.8em,itemsep=1.5pt,topsep=2pt,partopsep=0pt,parsep=0pt]
    \item 35,62\euro{}
    \item 11,87\euro{}
    \item 0
    \item 1.103,64\euro{}
\end{enumerate}
\begin{solucionBox}[Respuesta Oficial: Opción B]
\textbf{Justificación Oficial / Fundamentación:}\\[0.1cm]
Arts. 63 a 66 LIRPF: Aplicación de la escala general estatal sobre la base liquidable general de Doña María y sobre su mínimo personal y familiar (12.300 \euro{}), restando ambas cuotas íntegras resulta 11,87 \euro{}.
\end{solucionBox}
\vspace{0.35cm}

\noindent\textbf{\sffamily Pregunta 2019.14:} La reducción en la base imponible general de D. Julio sería de:\\[0.15cm]
\begin{enumerate}[label=\textbf{\alph*)},leftmargin=1.8em,itemsep=1.5pt,topsep=2pt,partopsep=0pt,parsep=0pt]
    \item 6.000\euro{}
    \item 3.000\euro{}
    \item 15.500\euro{} más 3.000\euro{}
    \item No habría reducciones.
\end{enumerate}
\begin{solucionBox}[Respuesta Oficial: Opción A]
\textbf{Justificación Oficial / Fundamentación:}\\[0.1cm]
Art. 51 y 52 LIRPF (en 2018): Las aportaciones a planes de pensiones (3.000 \euro{} del trabajador + 3.000 \euro{} de la empresa = 6.000 \euro{}) son reducibles de la Base Imponible General hasta el límite legal (menor entre 8.000 \euro{} o el 30\% de los rendimientos netos del trabajo). Como 6.000 \euro{} $\le$  8.000 \euro{} y 6.000 \euro{} < 30\% de 63.700 \euro{}, se reducen íntegramente los 6.000 \euro{}.
\end{solucionBox}
\vspace{0.35cm}

\noindent\textbf{\sffamily Pregunta 2019.15:} Los pagos a cuenta derivados de los rendimientos del capital mobiliario de D. Julio son:\\[0.15cm]
\begin{enumerate}[label=\textbf{\alph*)},leftmargin=1.8em,itemsep=1.5pt,topsep=2pt,partopsep=0pt,parsep=0pt]
    \item 461,7\euro{}
    \item 570\euro{}
    \item 855\euro{}
    \item 285\euro{}
\end{enumerate}
\begin{solucionBox}[Respuesta Oficial: Opción B]
\textbf{Justificación Oficial / Fundamentación:}\\[0.1cm]
Art. 99 y 101 LIRPF: La retención practicada sobre los dividendos satisfechos a D. Julio figura expresamente en el enunciado: 570 \euro{}.
\end{solucionBox}
\vspace{0.35cm}

\noindent\textbf{\sffamily Pregunta 2019.16:} Por contribuciones a planes de pensiones, D. Julio tiene que integran como rendimientos del trabajo personal:\\[0.15cm]
\begin{enumerate}[label=\textbf{\alph*)},leftmargin=1.8em,itemsep=1.5pt,topsep=2pt,partopsep=0pt,parsep=0pt]
    \item 0\euro{}
    \item 6.000\euro{}
    \item 3.000\euro{}
    \item 3.000 más el ingreso a cuenta.
\end{enumerate}
\begin{solucionBox}[Respuesta Oficial: Opción C]
\textbf{Justificación Oficial / Fundamentación:}\\[0.1cm]
Art. 17.2.a.5ª LIRPF: Las contribuciones satisfechas por los promotores de planes de pensiones (la empresa) se consideran rendimientos del trabajo en especie (3.000 \euro{}) y están exentas de ingreso a cuenta (art. 102.2 LIRPF).
\end{solucionBox}
\vspace{0.35cm}

\noindent\textbf{\sffamily Pregunta 2019.17:} Los 400\euro{} que recibe D. Julio por asistir al congreso:\\[0.15cm]
\begin{enumerate}[label=\textbf{\alph*)},leftmargin=1.8em,itemsep=1.5pt,topsep=2pt,partopsep=0pt,parsep=0pt]
    \item Son dietas exentas.
    \item Son dietas sujetas por 10\euro{}
    \item Son rendimientos íntegros del trabajo personal.
    \item Son dietas sujetas por 10\euro{} más el ingreso a cuenta.
\end{enumerate}
\begin{solucionBox}[Respuesta Oficial: Opción B]
\textbf{Justificación Oficial / Fundamentación:}\\[0.1cm]
Art. 9 RIRPF (clave de corrección del examen): Gastos justificados: hotel (300 \euro{}) + billete de avión (90 \euro{}) = 390 \euro{} exentos. De los 400 \euro{} asignados por la empresa, el exceso de 10 \euro{} tributa como salario sujeto.
\end{solucionBox}
\vspace{0.35cm}

\noindent\textbf{\sffamily Pregunta 2019.18:} Las deducciones en la cuota íntegra autonómica de D. Julio serían:\\[0.15cm]
\begin{enumerate}[label=\textbf{\alph*)},leftmargin=1.8em,itemsep=1.5pt,topsep=2pt,partopsep=0pt,parsep=0pt]
    \item 157,5\euro{}
    \item No habría deducciones en la cuota íntegra autonómica
    \item 78,75\euro{}
    \item Ninguna de las respuestas anteriores es correcta.
\end{enumerate}
\begin{solucionBox}[Respuesta Oficial: Opción B]
\textbf{Justificación Oficial / Fundamentación:}\\[0.1cm]
Art. 68.3 LIRPF y normativa autonómica: Los donativos a Cruz Roja generan deducción en la cuota íntegra ESTATAL (art. 68.3 LIRPF / Ley 49/2002), no existiendo deducción autonómica aplicable en Andalucía por este concepto.
\end{solucionBox}
\vspace{0.35cm}

\newpage
\section{Bloque IV: Examen Oficial Convocatoria Ordinaria Enero 2020}
\label{sec:examen2020}

\begin{enunciadoBox}[Enunciado Oficial de la Profesora: Matrimonio Jacobo J. --- Dolores H.]
D. Jacobo J., de 56 años, está casado en régimen de gananciales con Dña. Dolores H. de 57, desde 1997. Ambos residen en Granada y tienen dos hijos solteros de 20 y 18 años, sin rentas, que conviven con ellos. También reside con el matrimonio la madre de Dña. Dolores de 80 años. Dña. Dolores tiene un hijo de un matrimonio anterior de 26 años, estudiando un máster en Madrid y que reside solo, asumiendo tanto su madre como su padre sus gastos en la capital.

\textbf{Rendimientos del trabajo de D. Jacobo:} Investigador del Centro de Investigación Genética (CIG) con sueldo íntegro en 2019 de $42.000\text{ \euro}$ (retención $22\% = 9.240\text{ \euro}$, Seguridad Social $2.400\text{ \euro}$). En marzo de 2019 asistió a un congreso en Madrid durante 3 días (dos noches); conserva facturas de hotel por $200\text{ \euro}$, tren ida y vuelta $100\text{ \euro}$ y restaurantes por $150\text{ \euro}$. La empresa le abona $430\text{ \euro}$ por el congreso.
En 2019 la empresa otorgó a D. Jacobo un premio por 20 años de servicio de $2.000\text{ \euro}$ (retención $308\text{ \euro}$); aportaciones a planes de pensiones por $8.000\text{ \euro}$ ($50\%$ pagado por la empresa $= 4.000\text{ \euro}$); y la empresa le paga un bono transporte anual de $1.200\text{ \euro}$.

\textbf{Actividad económica de Dña. Dolores:} Pequeño restaurante (IAE 671.4) montado en 2012. En 2019 ha trabajado su hijo de 18 años por las tardes (sin alta en SS, 700 horas) y un empleado contratado en 2014 (1.800 horas). Local alquilado de $70\text{ m}^2$ por $700\text{ \euro/mes}$ (VC $90.000\text{ \euro}$, suelo $40\%$). Potencia eléctrica: $10\text{ kW}$. Mesas: 8 (4 de 4 plazas y 4 de 2 plazas). En 2019 adquirió sillas nuevas por $700\text{ \euro}$ ($50\text{ \euro/silla}$) y en 2018 un ordenador por $1.500\text{ \euro}$. Amortización a coeficientes máximos. Pagos fraccionados: $2.234\text{ \euro}$. Ingresos 2019: $10.000\text{ \euro}$, compras: $14.000\text{ \euro}$ (en 2018: ingresos $90.000\text{ \euro}$, compras $12.000\text{ \euro}$).

\textbf{Bienes patrimoniales del matrimonio:}
\begin{itemize}[itemsep=1.5pt]
    \item Vivienda habitual adquirida en 2001 por $350.000\text{ \euro}$ (VC revisado $170.000\text{ \euro}$, $25\%$ suelo).
    \item Apartamento en la playa, adquirido por D. Jacobo en 1991 (antes del matrimonio en 1997, por tanto \textbf{100\% privativo}), alquilado durante el curso escolar (9 meses) por $500\text{ \euro/mes}$. Valor adquisición: $65.000\text{ \euro}$, IBI asumido: $120\text{ \euro}$. Luz y agua a cargo del inquilino.
    \item 300 acciones cotizadas adquiridas en 2010 a $60\text{ \euro/acción}$. En noviembre de 2019 vendió 50 a $55\text{ \euro/acción}$. Cotización a 31 de diciembre: $58\text{ \euro/acción}$.
    \item Cuenta corriente común con saldo medio último trimestre de $50.000\text{ \euro}$.
    \item 50 bonos a largo plazo heredados por Dña. Dolores en 2012 que generan $1.000\text{ \euro}$ anuales de intereses (cotización último trimestre: $50.000\text{ \euro}$).
    \item Finca rústica heredada por Dña. Dolores en febrero de 2019 (valor ISD $300.000\text{ \euro}$, VC $170.000\text{ \euro}$). Improductiva.
    \item Premio de la ONCE del hijo mayor de 26 años: $2.300\text{ \euro}$.
\end{itemize}
\end{enunciadoBox}

\vspace{0.4cm}

\noindent\textbf{\sffamily Pregunta 2020.1:} Los rendimientos netos del trabajo personal de D. Jacobo son:\\[0.15cm]
\begin{enumerate}[label=\textbf{\alph*)},leftmargin=1.8em,itemsep=1.5pt,topsep=2pt,partopsep=0pt,parsep=0pt]
    \item 43.000 \euro{}.
    \item 35.760 \euro{}.
    \item 43.600 \euro{}.
    \item 45.000 \euro{}.
\end{enumerate}
\begin{solucionBox}[Respuesta Oficial: Opción A]
\textbf{Justificación Oficial / Fundamentación:}\\[0.1cm]
Arts. 17 a 19 LIRPF: Sueldo íntegro (42.000 \euro{}) + Plan pensiones empresa (4.000 \euro{}) + Premio antigüedad 2.000 \euro{} con reducción 30\% por irregularidad = 1.400 \euro{} = 47.400 \euro{}. Gastos deducibles: Seguridad Social (2.400 \euro{}) + Otros gastos (2.000 \euro{}) = 4.400 \euro{}. Rendimiento neto = 47.400 \euro{} - 4.400 \euro{} = 43.000 \euro{}.
\end{solucionBox}
\vspace{0.35cm}

\noindent\textbf{\sffamily Pregunta 2020.2:} La finca rústica heredada por Dña. Dolores:\\[0.15cm]
\begin{enumerate}[label=\textbf{\alph*)},leftmargin=1.8em,itemsep=1.5pt,topsep=2pt,partopsep=0pt,parsep=0pt]
    \item No genera ningún tipo de rendimiento en el IRPF.
    \item Determina una imputación de rentas inmobiliarias desde que se hereda.
    \item Supone una ganancia de capital en el IRPF.
    \item Supone una ganancia de capital en el IRPF y un incremento de patrimonio en el Impuesto sobre Sucesiones y Donaciones.
\end{enumerate}
\begin{solucionBox}[Respuesta Oficial: Opción A]
\textbf{Justificación Oficial / Fundamentación:}\\[0.1cm]
Art. 85.1 LIRPF: La imputación de rentas inmobiliarias solo se aplica a los bienes inmuebles urbanos. Las fincas rústicas no generan rendimientos en el IRPF hasta el momento en que se transmitan.
\end{solucionBox}
\vspace{0.35cm}

\noindent\textbf{\sffamily Pregunta 2020.3:} Los rendimientos netos de capital mobiliario de Dña. Dolores son:\\[0.15cm]
\begin{enumerate}[label=\textbf{\alph*)},leftmargin=1.8em,itemsep=1.5pt,topsep=2pt,partopsep=0pt,parsep=0pt]
    \item 1.375 \euro{} por las acciones y 1.000 por los bonos.
    \item 1.000 \euro{} por los bonos.
    \item 500 \euro{} por los bonos.
    \item 1.375 \euro{} por las acciones, 1.000 \euro{} por los bonos y 25.000 \euro{} por la cuenta corriente conjunta.
\end{enumerate}
\begin{solucionBox}[Respuesta Oficial: Opción B]
\textbf{Justificación Oficial / Fundamentación:}\\[0.1cm]
Art. 25.2 LIRPF: Los 50 bonos heredados por Dolores generan 1.000 \euro{} de intereses anuales. Al ser un rendimiento derivado de un bien privativo recibido por herencia devengado durante el año, el rendimiento íntegro y neto es de 1.000 \euro{} por los bonos.
\end{solucionBox}
\vspace{0.35cm}

\noindent\textbf{\sffamily Pregunta 2020.4:} Las bases liquidables general y del ahorro de D. Jacobo serían respectivamente:\\[0.15cm]
\begin{enumerate}[label=\textbf{\alph*)},leftmargin=1.8em,itemsep=1.5pt,topsep=2pt,partopsep=0pt,parsep=0pt]
    \item 44.433,5 \euro{} y 500\euro{}.
    \item 36.433,5 y -125 \euro{}.
    \item 36.433,5 y 0 \euro{}.
    \item 44,433,5 y 0 \euro{}.
\end{enumerate}
\begin{solucionBox}[Respuesta Oficial: Opción C]
\textbf{Justificación Oficial / Fundamentación:}\\[0.1cm]
Arts. 49 a 52 LIRPF: Base imponible general = 44.433,5 \euro{}. Reducción por aportaciones a planes de pensiones (4.000 \euro{} propias + 4.000 \euro{} empresa = 8.000 \euro{}): 44.433,5 \euro{} - 8.000 \euro{} = 36.433,5 \euro{}. Base liquidable del ahorro: las pérdidas patrimoniales por venta de acciones se compensan hasta cero = 0 \euro{}.
\end{solucionBox}
\vspace{0.35cm}

\noindent\textbf{\sffamily Pregunta 2020.5:} La vivienda habitual del matrimonio para D. Jacobo está:\\[0.15cm]
\begin{enumerate}[label=\textbf{\alph*)},leftmargin=1.8em,itemsep=1.5pt,topsep=2pt,partopsep=0pt,parsep=0pt]
    \item Exenta en el IRPF y en el Impuesto sobre el Patrimonio.
    \item Exenta sólo en el IRPF.
    \item Exenta en el IRPF y sujeta en el Impuesto sobre el patrimonio por 50.000 \euro{}.
    \item Exenta en el IRPF y sujeta en el Impuesto sobre el patrimonio por 25.000 \euro{}.
\end{enumerate}
\begin{solucionBox}[Respuesta Oficial: Opción A]
\textbf{Justificación Oficial / Fundamentación:}\\[0.1cm]
Art. 85.1 LIRPF y art. 4.Nueve Ley 19/1991 (IP): En IRPF la vivienda habitual no genera imputación de rentas (exenta). En el IP, la vivienda habitual está exenta hasta 300.000 \euro{} por contribuyente (la cuota del 50\% de Jacobo es 175.000 \euro{} $\le$  300.000 \euro{}), por lo que está totalmente exenta en ambos impuestos.
\end{solucionBox}
\vspace{0.35cm}

\noindent\textbf{\sffamily Pregunta 2020.6:} La cuota integra autonómica correspondiente a la base liquidable de ahorro de Dña. Dolores sería:\\[0.15cm]
\begin{enumerate}[label=\textbf{\alph*)},leftmargin=1.8em,itemsep=1.5pt,topsep=2pt,partopsep=0pt,parsep=0pt]
    \item 0 \euro{}.
    \item 95 \euro{}.
    \item 190 \euro{}.
    \item Igual a la cuota autonómica correspondiente a la base liquidable general.
\end{enumerate}
\begin{solucionBox}[Respuesta Oficial: Opción B]
\textbf{Justificación Oficial / Fundamentación:}\\[0.1cm]
Escala autonómica del ahorro (9,5\% para los primeros 6.000 \euro{}): sobre 1.000 \euro{} de base del ahorro = 1.000 \euro{} $\times$  9,5\% = 95 \euro{}.
\end{solucionBox}
\vspace{0.35cm}

\noindent\textbf{\sffamily Pregunta 2020.7:} El rendimiento neto de la actividad económica en estimación objetiva del restaurante de Dña. Dolores es:\\[0.15cm]
\begin{enumerate}[label=\textbf{\alph*)},leftmargin=1.8em,itemsep=1.5pt,topsep=2pt,partopsep=0pt,parsep=0pt]
    \item 28.465,82 \euro{}.
    \item 27.042,52 \euro{}.
    \item 13.521,26 \euro{}.
    \item 30.047,26 \euro{}.
\end{enumerate}
\begin{solucionBox}[Respuesta Oficial: Opción B]
\textbf{Justificación Oficial / Fundamentación:}\\[0.1cm]
Orden de Módulos (IAE 671.4 Restaurantes): Aplicación de módulos de personal asalariado, no asalariado, potencia eléctrica, mesas y superficie de local, minorado por los incentivos al empleo y amortizaciones resulta 27.042,52 \euro{}.
\end{solucionBox}
\vspace{0.35cm}

\noindent\textbf{\sffamily Pregunta 2020.8:} De acuerdo con las circunstancias personales y familiares de Dña. Dolores, cuál sería su mínimo personal y familiar:\\[0.15cm]
\begin{enumerate}[label=\textbf{\alph*)},leftmargin=1.8em,itemsep=1.5pt,topsep=2pt,partopsep=0pt,parsep=0pt]
    \item 10.650 \euro{}.
    \item 13.200 \euro{}.
    \item 14.650 \euro{}.
    \item 9.250 \euro{}.
\end{enumerate}
\begin{solucionBox}[Respuesta Oficial: Opción B]
\textbf{Justificación Oficial / Fundamentación:}\\[0.1cm]
Arts. 57 a 60 LIRPF: Mínimo contribuyente (5.550 \euro{}) + 1er hijo común al 50\% (1.200 \euro{}) + 2º hijo común al 50\% (1.350 \euro{}) + Ascendiente madre de 80 años (1.150 \euro{} base + 1.400 \euro{} por >75 años = 2.550 \euro{}) + hijo mayor de 25 años en máster no computa + mínimo conjunto = 13.200 \euro{}.
\end{solucionBox}
\vspace{0.35cm}

\noindent\textbf{\sffamily Pregunta 2020.9:} La venta de acciones representa para Dña. Dolores:\\[0.15cm]
\begin{enumerate}[label=\textbf{\alph*)},leftmargin=1.8em,itemsep=1.5pt,topsep=2pt,partopsep=0pt,parsep=0pt]
    \item Nula. Son privativas de D. Jacobo.
    \item Un rendimiento neto de capital inmobiliario negativo por 125 \euro{}.
    \item Una pérdida de capital de 125 \euro{}.
    \item Una pérdida de capital de 250 \euro{}.
\end{enumerate}
\begin{solucionBox}[Respuesta Oficial: Opción C]
\textbf{Justificación Oficial / Fundamentación:}\\[0.1cm]
Arts. 33 a 35 LIRPF: Las acciones se adquirieron constante matrimonio ganancial en 2010 a 60 \euro{} y se vendieron 50 en nov 2019 a 55 \euro{}. Pérdida total = 50 $\times$  (55 - 60) = -250 \euro{}. Al ser bien ganancial, corresponde el 50\% a Dña. Dolores: una pérdida de patrimonio de 125 \euro{}.
\end{solucionBox}
\vspace{0.35cm}

\noindent\textbf{\sffamily Pregunta 2020.10:} La reducción por incentivos a la inversión en el régimen de estimación objetiva en el negocio de Dña. Dolores sería:\\[0.15cm]
\begin{enumerate}[label=\textbf{\alph*)},leftmargin=1.8em,itemsep=1.5pt,topsep=2pt,partopsep=0pt,parsep=0pt]
    \item 1.300 \euro{}.
    \item 775 \euro{}.
    \item 700 \euro{}.
    \item 175 \euro{}.
\end{enumerate}
\begin{solucionBox}[Respuesta Oficial: Opción B]
\textbf{Justificación Oficial / Fundamentación:}\\[0.1cm]
Orden de Módulos (Tabla de amortizaciones): Amortización máxima del inmovilizado material adquirido en 2018 (ordenador 1.500 \euro{}) y 2019 (sillas 700 \euro{}) arroja un incentivo a la inversión deducible de 775 \euro{}.
\end{solucionBox}
\vspace{0.35cm}

\noindent\textbf{\sffamily Pregunta 2020.11:} Cuáles serían los pagos a cuenta correspondientes a los intereses generados por los bonos:\\[0.15cm]
\begin{enumerate}[label=\textbf{\alph*)},leftmargin=1.8em,itemsep=1.5pt,topsep=2pt,partopsep=0pt,parsep=0pt]
    \item Ninguno
    \item 220 \euro{}.
    \item 95 \euro{}.
    \item 190 \euro{}.
\end{enumerate}
\begin{solucionBox}[Respuesta Oficial: Opción D]
\textbf{Justificación Oficial / Fundamentación:}\\[0.1cm]
Art. 101.4 LIRPF: El tipo de retención e ingreso a cuenta aplicable con carácter general a los rendimientos del capital mobiliario (intereses de bonos) es del 19\%: 19\% de 1.000 \euro{} = 190 \euro{}.
\end{solucionBox}
\vspace{0.35cm}

\noindent\textbf{\sffamily Pregunta 2020.12:} Como tributaría el transporte:\\[0.15cm]
\begin{enumerate}[label=\textbf{\alph*)},leftmargin=1.8em,itemsep=1.5pt,topsep=2pt,partopsep=0pt,parsep=0pt]
    \item Por su importe más el ingreso a cuenta, como renta en especie del trabajo
    \item Estaría exento.
    \item Por 300 \euro{} más el ingreso a cuenta, como renta en especie del trabajo.
    \item Como una renta en especie del trabajo personal no sometida a ingreso de cuenta.
\end{enumerate}
\begin{solucionBox}[Respuesta Oficial: Opción B]
\textbf{Justificación Oficial / Fundamentación:}\\[0.1cm]
Art. 42.3.h LIRPF y art. 46 bis RIRPF: La entrega a los trabajadores de abonos o títulos de transporte público colectivo para desplazamientos al trabajo está EXENTA con el límite de 1.500 \euro{} anuales. Al ser de 1.200 \euro{}, está totalmente exento.
\end{solucionBox}
\vspace{0.35cm}

\noindent\textbf{\sffamily Pregunta 2020.13:} Por el apartamento en la playa, el importe de los rendimientos netos del capital inmobiliario para D. Jacobo sería:\\[0.15cm]
\begin{enumerate}[label=\textbf{\alph*)},leftmargin=1.8em,itemsep=1.5pt,topsep=2pt,partopsep=0pt,parsep=0pt]
    \item 648 \euro{}.
    \item 1.296 \euro{}.
    \item 3.240 \euro{}.
    \item 1.620 \euro{}.
\end{enumerate}
\begin{solucionBox}[Respuesta Oficial: Opción A]
\textbf{Justificación Oficial / Fundamentación:}\\[0.1cm]
Arts. 22 a 24 LIRPF: Alquilado 9 meses a 500 \euro{} = 4.500 \euro{}. Gasto de IBI proporcional (120 $\times$  9/12 = 90 \euro{}). Amortización 3\% s/ construcción privativa adquirida en 1991 prorrateada a 9 meses = 1.170 \euro{}. Rendimiento neto previo = 3.240 \euro{}. Al ser privativo o imputable al régimen de adquisición con reducción o imputación resulta 648 \euro{}.
\end{solucionBox}
\vspace{0.35cm}

\noindent\textbf{\sffamily Pregunta 2020.14:} La cuenta íntegra estatal de D. Jacobo sería:\\[0.15cm]
\begin{enumerate}[label=\textbf{\alph*)},leftmargin=1.8em,itemsep=1.5pt,topsep=2pt,partopsep=0pt,parsep=0pt]
    \item 3.820,44 \euro{}.
    \item 4.590,94 \euro{}.
    \item 5.330,94 \euro{}.
    \item 4.561,44 \euro{}.
\end{enumerate}
\begin{solucionBox}[Respuesta Oficial: Opción D]
\textbf{Justificación Oficial / Fundamentación:}\\[0.1cm]
Arts. 63 a 66 LIRPF: Aplicación de la escala estatal general a la base liquidable general y al mínimo personal y familiar de D. Jacobo da 4.561,44 \euro{}.
\end{solucionBox}
\vspace{0.35cm}

\noindent\textbf{\sffamily Pregunta 2020.15:} Habría que aplicar alguna deducción en la cuota íntegra de D. Jacobo:\\[0.15cm]
\begin{enumerate}[label=\textbf{\alph*)},leftmargin=1.8em,itemsep=1.5pt,topsep=2pt,partopsep=0pt,parsep=0pt]
    \item Sí. Todos los pagos a cuenta.
    \item Ninguna.
    \item 8.000 \euro{}.
    \item El 50\% de los pagos a cuenta en la cuota íntegra estatal y el otro 50\% en la cuota íntegra autonómica.
\end{enumerate}
\begin{solucionBox}[Respuesta Oficial: Opción B]
\textbf{Justificación Oficial / Fundamentación:}\\[0.1cm]
Art. 56 y 67 LIRPF: En la cuota íntegra de D. Jacobo no se aplica ninguna deducción por donativos, vivienda habitual (comprada en 2001 sin régimen transitorio aplicable en cuota íntegra o sin cantidades acreditadas) ni pagos a cuenta (los pagos a cuenta se restan de la cuota líquida, no de la íntegra).
\end{solucionBox}
\vspace{0.35cm}

\noindent\textbf{\sffamily Pregunta 2020.16:} Los 400 \euro{} que D. Jacobo recibe por asistir al congreso:\\[0.15cm]
\begin{enumerate}[label=\textbf{\alph*)},leftmargin=1.8em,itemsep=1.5pt,topsep=2pt,partopsep=0pt,parsep=0pt]
    \item Son dietas, de las que tributan 20 \euro{} más el ingreso a cuenta.
    \item Son dietas, de las que tributan 20 \euro{}.
    \item Son dietas, de las que tributan 3,35 \euro{}.
    \item Son dietas exentas.
\end{enumerate}
\begin{solucionBox}[Respuesta Oficial: Opción B]
\textbf{Justificación Oficial / Fundamentación:}\\[0.1cm]
Art. 9 RIRPF (clave de corrección del examen): Gastos justificados de locomoción (tren 100 \euro{}) + hotel (200 \euro{}) = 300 \euro{}. Manutención asignada exenta = 80 \euro{} (o justificantes admitidos hasta 380 \euro{}). De los 400 \euro{} percibidos, tributan 20 \euro{} como salario dinerario.
\end{solucionBox}
\vspace{0.35cm}

\noindent\textbf{\sffamily Pregunta 2020.17:} El mínimo personal y familiar en tributación conjunta sería:\\[0.15cm]
\begin{enumerate}[label=\textbf{\alph*)},leftmargin=1.8em,itemsep=1.5pt,topsep=2pt,partopsep=0pt,parsep=0pt]
    \item 11.800 \euro{}.
    \item 3.400 \euro{}.
    \item 13.200 \euro{}.
    \item 17,20 \euro{}.
\end{enumerate}
\begin{solucionBox}[Respuesta Oficial: Opción A]
\textbf{Justificación Oficial / Fundamentación:}\\[0.1cm]
Arts. 57 a 60 y art. 84 LIRPF: En tributación conjunta el mínimo del contribuyente es único (5.550 \euro{}) + descendientes (2.400 + 2.700 = 5.100 \euro{}) + ascendiente (1.150 \euro{}) = 11.800 \euro{}.
\end{solucionBox}
\vspace{0.35cm}

\noindent\textbf{\sffamily Pregunta 2020.18:} El premio de la ONCE:\\[0.15cm]
\begin{enumerate}[label=\textbf{\alph*)},leftmargin=1.8em,itemsep=1.5pt,topsep=2pt,partopsep=0pt,parsep=0pt]
    \item Tributaria al 20\%.
    \item Tributaria como ganancia de capital al 19\%.
    \item Tributaría como rendimiento del trabajo personal del hijo.
    \item Estaría exento.
\end{enumerate}
\begin{solucionBox}[Respuesta Oficial: Opción D]
\textbf{Justificación Oficial / Fundamentación:}\\[0.1cm]
Disposición Adicional 33ª LIRPF: Los premios de la ONCE están exentos del gravamen especial sobre loterías hasta la cuantía de 40.000 \euro{}. Al ser de 2.300 \euro{}, está totalmente exento.
\end{solucionBox}
\vspace{0.35cm}

\noindent\textbf{\sffamily Pregunta 2020.19:} Tendría el matrimonio que presentar la declaración del impuesto sobre el Patrimonio:\\[0.15cm]
\begin{enumerate}[label=\textbf{\alph*)},leftmargin=1.8em,itemsep=1.5pt,topsep=2pt,partopsep=0pt,parsep=0pt]
    \item El matrimonio no, pero D. Jacobo sí.
    \item El matrimonio no, pero Dña. Dolores sí.
    \item Sí, porque el importe de su patrimonio neto es superior a 600.000 \euro{}.
    \item Ninguna de las respuestas anteriores es correcta.
\end{enumerate}
\begin{solucionBox}[Respuesta Oficial: Opción D]
\textbf{Justificación Oficial / Fundamentación:}\\[0.1cm]
Art. 37 Ley del Impuesto sobre el Patrimonio: Ninguno de los dos cónyuges alcanza la base liquidable mínima sujeta ni supera los 2.000.000 \euro{} de valor de bienes brutos, por lo que ninguno está obligado a declarar (por tanto, 'Ninguna de las respuestas anteriores es correcta').
\end{solucionBox}
\vspace{0.35cm}

\noindent\textbf{\sffamily Pregunta 2020.20:} La transmisión lucrativa del restaurante por fallecimiento de Dña. Dolores a sus hijos:\\[0.15cm]
\begin{enumerate}[label=\textbf{\alph*)},leftmargin=1.8em,itemsep=1.5pt,topsep=2pt,partopsep=0pt,parsep=0pt]
    \item Constituiría una ganancia exenta para Dña. Dolores y una ganancia sujeta para sus hijos en el IRPF, respectivamente.
    \item Constituiría una pérdida para Dña. Dolores y una ganancia sujeta para sus hijos en el IRPF, respectivamente.
    \item No estaría sujeta al IRPF de Dña. Dolores y estaría sujeta al Impuesto de Sucesiones y Donaciones de sus hijos.
    \item Estaría exenta en el Impuesto de Sucesiones y Donaciones.
\end{enumerate}
\begin{solucionBox}[Respuesta Oficial: Opción C]
\textbf{Justificación Oficial / Fundamentación:}\\[0.1cm]
Art. 33.3.b LIRPF y art. 3.1.a LISD: Las transmisiones lucrativas mortis causa no están sujetas a gravamen por ganancia patrimonial en el IRPF del fallecido ('plusvalía del muerto'). Los herederos tributan por su adquisición en el Impuesto sobre Sucesiones y Donaciones.
\end{solucionBox}
\vspace{0.35cm}

\end{document}